% !PRINTED 294
\weberpage{294}
Wir haben nun im §.\,154, 7. des ersten Bandes den Satz bewiesen, dass der Grad
einer transitiven Permutationsgruppe immer durch die Anzahl der permutirten
Ziffern theilbar ist. Hier ist aber $f_1(x)$ irreducibel und demnach die Gruppe
$Q$ transitiv. Es ist also $q$ durch den Grad von $f_1(x)$, d.\,h. durch $m_1$
theilbar, und folglich ist nach (4) auch $p_\lambda$ durch $m_1$ theilbar.

Den Grad $p_{\lambda-1}$ von $P_{\lambda-1}$ erhalten wir nach der Bedeutung von
$j_\lambda$ in der Form:
\weq{5}{p_{\lambda-1} = j_\lambda\, p_\lambda.}
Nun ist $p_\lambda$ ein Vielfaches von $m_1$, $j_\lambda$ ein Vielfaches von $h$,
und $h\, m_1 = m$ gleich dem Grade von $f(x)$. Demnach ist
\weq{6}{p_{\lambda-1} = k\, m}
ein Vielfaches von $m$.

Aus der Bedeutung der $j$ ergiebt sich aber [§.\,73, (2)]
$p_{\lambda-1} = j_\lambda j_{\lambda-1} \ldots j_\mu$, % [sic] second factor $j_{\lambda-1}$ as printed; equation (7) gives $j_{\lambda+1}$.
und folglich haben wir die
Relation
\weq{7}{j_\lambda j_{\lambda+1} j_{\lambda+2} \ldots j_\mu = k\, m,}
woraus sich eine sehr merkwürdige Folgerung ziehen lässt.

Angenommen, es werde die irreducible Function $f(x)$ durch successive Adjunction
von Wurzeln cyklischer Gleichungen reducibel, dann lässt sich nach §.\,176, Bd. I
die Compositionsreihe (1) so geordnet annehmen, dass die Indices
$j_1, j_2, \ldots, j_\lambda$ Primzahlen sind. Es ist also, da $h$ ein Theiler von
$j_\lambda$ ist, $j_\lambda = h$, und $j_\lambda$ ist nach (3) eine in $m$
aufgehende Primzahl.

Wenn nun in $m$ ausser $j_\lambda$ noch eine zweite Primzahl $p$ aufgeht, so muss
nach (7) einer der Factoren $j_\lambda, j_{\lambda+1}, j_{\lambda+2}, \ldots, j_\mu$
durch $p$ theilbar sein, und sie können also gewiss nicht alle gleich $j_\lambda$
sein.

Daraus aber folgt nach dem Satze III, §.\,7, dass man eine Compositionsreihe von
$P_\lambda$:
\weq{8}{P_\lambda,\ P_{\lambda+1},\ \ldots,\ P_{\mu-1},\ 1}
so finden kann, dass unter den Gruppen $P_\lambda, P_{\lambda+1}, \ldots,
P_{\mu-1}$ eine, etwa $P_\nu$, ein \gesperrt{Normaltheiler von} $P$ ist.

Dieses $P_\nu$ ist aber als Theiler der intransitiven Gruppe $P_\lambda$ selbst
intransitiv, und daher muss $P$ selbst \gesperrt{imprimitiv} sein (Bd. I, §.\,158, 2).

Wir sprechen dies in folgender Form als Satz aus:

\gesperrt{1. Wenn im Grade einer irreduciblen Gleichung mehrere verschiedene
Primzahlen aufgehen, so}
%CONT>>
